\section{Výběr základních parametrů}

Tato kapitola se bude zabývat základními výpočty vlastností robota.

\subsection{Rozměr}
Prvním omezením je velikost samotného \ac{dps}. Jestliže bude využita sleva u \ac{jlc}, která platí pro maximální rozměr 100x100mm. Další limitací jsou pouze dvě vrstvy, které mohou být využity. Při překročení některé z těchto limitací nevychází jedno \ac{dps} na 5\$, ale na řádově šestinásobek.

Na 3D tiskárně byla vytištěna plastová destička o rozměrech 100x100x1.6mm, pro lepší vizualizaci rozměru \ac{dps}.

\subsection{Motory}
U těch je hlavní točivý moment, který musí odpovívat rovnici\\
\begin{equation}\label{eq:vypocet sily motoru}
	F = m \cdot q \cdot l \cdot \sin \left( \theta \right) \cdot \frac{S}{2}
\end{equation}
\newline
Význam proměnných:\\
m = hmotnost\\
q = gravitační zrychlení země [$9.81m \cdot s^{-2}$]\\
l = výška robota\\
$\theta$ = úhel naklonění robota vůči nulové ose\\
S = safety factor [2]\\
2 = počet motorů (stačí pouze polovina síly potřebné pro zvrácení robota)
\newline

Maximální hmotnost robota bude 300g.\\
maximální výška je 150mm. Jedná se o maximální výšku PCB (100mm) + maximální poloměr kol (50mm).
Zrychlení se dá vypočítat podle vztahu $a = q \cdot sin(\theta)$\\
Pro výpočtu maximálního nákllonu můžeme využít vztah $\theta = sin^{-1}(\frac{a}{q})$\\
Pro praktické zrychlení alespoň $1m*s^{-2}$ je potřeba využít úhel $\theta = 5.85°$\\
Po dosazení do rovnice \ref{eq:vypocet sily motoru} vyjde požadovaný točivý moment $45mNm$\\
Vztah pro přepočet $Nm$ na $kg \cdot cm$ je $kg \cdot cm = \frac{Nm}{q} \ cdot 100$\\
Po dosazení vyjde točivý moment jednoho motoru $0.459kg \cdot cm$
\begin{figure}[htbp!]
	\includegraphics[width = \textwidth]{motor_table.png}
	\label{table_with_motor_parameters}
	\caption{Tabulka obsahuje parametry mikromotorů 12GA-N20}
	\cite{aliexpress_mikromotor}, upraveno
\end{figure}

Z tabulky \ref{table_with_motor_parameters}, parametry mikromotorů 12GA-N20 při napětí 12V vyhovuje motor o rychlosti $1000rpm$

\subsection{Baterie}
Jeden článek baterie 18650 má hmotnost 46g.

\begin{table}[htbp!]
	\centering
	\begin{tabularx}{\textwidth}{|X|X|X|X|X|}\hline
		Počet článků & Přibližná $m$ $\left[g\right]$ & Průměr. $U$ $\left[ V \right]$ & Max. $P$ $\left[ W \right]$ \\ \hline
		1S           & 46 + 15     & 3.6                            & 36                          \\ \hline
		2S           & 92 + 25     & 7.2                            & 72                          \\ \hline
		3S           & 138 + 35    & 11.8                           & 118                         \\ \hline
	\end{tabularx}
	\label{tab:tabulka 18650}
	\caption{Tabulka porovnání počtu článků 18650}
\end{table}
Při použití tří článků by sice nebyly potřebné step up moduly, avšak kvůli velikosti není kam tři články umístit.
maximální výkon byl počítán ze vztahu $p=U \cdot I$ kde i je průměrné napětí a I je 10A.
Hmotnost článků je známá, hmotnost držáku je však pouze odhadovaná.
\newline

\begin{table}[htbp!]
	\centering
	\begin{tabularx}{\textwidth}{|X|X|X|X|X|}\hline
		Název modulu & Přibliž. spotř. $P$ $\left[W\right]$ & Požad. $U$ $\left[ V \right]$ & Max. $P$ $\left[ W \right]$ \\ \hline
		1S           & 46 + 15     & 3.6                            & 36                          \\ \hline
		2S           & 92 + 25     & 7.2                            & 72                          \\ \hline
		3S           & 138 + 35    & 11.8                           & 118                         \\ \hline
	\end{tabularx}
	\label{tab:tabulka_prikonnu}
	\caption{Tabulka příkonů}
\end{table}

\subsection{Motory, 2. kolo }
Jelikož již známe přibližnou hmotabulka 18650

\textbf{Round robin}
Jedna hlavni funkce bude vše řídit.

\textbf{Roud robin s Interrupty}

jednoduchost

Pouze jedno jádro

\textbf{\ac{freertos}}
Jedna funkce přiřadí úkony. Operační systém se poté postará o přiřazení operací jednotlivým jádrům.

Možnost použití obou jader


\textbf{Kombinace}

Hlavní funkce přiřadí operaci s tím, že jednomu jádru přiřadí pouze jednu funkci a druhému jádru přiřadí ostatní funkce. Nestane se tedy, že by byl přerušen hlavní loop starající se o balancování robota a robot spadnul. Zároveň se mohou vykonávat asynchronní operace jako komunikace s okolím, či logování dat na SD kartu.

Základními

$F=m\cdot l\cdot\sin\left[ \alpha \right] \cdot\frac{S}{2}$ [kg/cm]
$m\cdot l\cdot q\cdot\sin\left(a\right)$ [Ncm]
$\alpha=\sin^{-1}\left(\frac{a}{mlq}\right)$
$a=10ms^{-2}$
$m=0.3kg$
$l=15$
$q=9.81$
$S=2$